\section{Условия обучения}
\label{sec:experiments}

\subsection{Гиперпараметры}

Все модели обучаются с едиными гиперпараметрами (таблица~\ref{tab:hparams}), чтобы различия в результатах можно было однозначно отнести к архитектуре, а не к условиям обучения.

\begin{table}[h]
\centering
\caption{Гиперпараметры обучения}
\label{tab:hparams}
\begin{tabular}{ll}
\toprule
\textbf{Параметр} & \textbf{Значение} \\
\midrule
Оптимизатор             & AdamW, weight decay $= 0.01$ \\
Learning rate           & $2 \times 10^{-5}$ (энкодер + голова) \\
Learning rate (CRF)     & $10^{-3}$ (только CRF-слой) \\
Планировщик             & Linear warmup (1 эпоха) + линейный decay \\
Batch size              & 16 (8 для scibert\_concat4) \\
Максимум эпох           & 10 \\
Early stopping patience & 3 (по dev macro F1) \\
Случайное зерно         & 42 \\
Смешанная точность      & AMP fp16 на GPU \\
\bottomrule
\end{tabular}
\end{table}

\textbf{Выбор числа эпох.} Лимит в 10 эпох --- верхняя граница, реальным критерием остановки служит early stopping: обучение прекращается, если dev macro F1 не улучшается 3 эпохи подряд, и сохраняется лучший чекпоинт. На датасете SciERC (~1500 обучающих предложений) BERT-scale модели эмпирически выходят на плато к 6--8 эпохе, а дальнейшее снижение train loss сопровождается ростом dev loss --- признак переобучения. Увеличение лимита до 15--20 эпох не изменило бы результатов: early stopping остановил бы большинство прогонов раньше. В экспериментах большинство моделей остановились на 7--10 эпохе, что подтверждает достаточность выбранного лимита (рис.~\ref{fig:curves_v1}).

\textbf{Выбор learning rate.} Значение $2 \times 10^{-5}$ --- стандартное для fine-tuning BERT-scale энкодеров~\cite{devlin2019bert}. Для CRF-слоя используется отдельный, более высокий lr $10^{-3}$: CRF не имеет предобученных весов и требует более быстрой начальной сходимости.

\subsection{Инфраструктура}

Обучение проводилось на платформе Kaggle на двух GPU NVIDIA T4 (по 15.6 GB памяти). Для параллелизма применялся \texttt{torch.nn.DataParallel}, за исключением \texttt{scibert\_concat4}: эта модель захватывает скрытые состояния промежуточных слоёв через forward-хуки, что несовместимо с DataParallel. Тестовая оценка производится на чекпоинте лучшей эпохи по dev F1.
